\documentclass[11pt]{amsart}
\usepackage{amssymb,mathtools}
\usepackage[margin=1.15in]{geometry}
\usepackage[hidelinks]{hyperref}
\newtheorem{theorem}{Theorem}[section]
\newtheorem{lemma}[theorem]{Lemma}
\theoremstyle{definition}\newtheorem{definition}[theorem]{Definition}
\newcommand{\UU}{\mathcal U}
\newcommand{\Hom}{\operatorname{Hom}}
\newcommand{\Spec}{\operatorname{Spec}}
\newcommand{\Nat}{\operatorname{Nat}}
\newcommand{\tr}[1]{\left\Vert#1\right\Vert}
\newcommand{\mv}{\operatorname{move}}
\title[Relative totalisation over a base ring]{Relative totalisation over a base ring\\A proof with explicit changes of structure map}
\date{October 29, 2026}
\begin{document}\raggedbottom\maketitle

\section{Tests and transport}
Fix a small commutative ring $k$ with set carrier. A test algebra is a
pair $(B,\iota)$ of a small set-carrier ring and a map $\iota:k\to B$.
An algebra map $(B,\iota)\to(C,\kappa)$ consists of $h:B\to C$ and
$h\iota=\kappa$. One may instead give this equation pointwise; ring-map
extensionality identifies the two forms. Both equations are propositions.
A relative functor $F$ assigns small sets to these algebras and functions
to algebra maps, satisfying the identity and composition laws.

For fixed $B$ write $F_B(\iota)=F(B,\iota)$. For a path
$p:\iota=\kappa$, let $\mv_F(p,x)$ denote transport in this family.
Transport at reflexivity is identity, by definition. Path induction gives
transport along successive paths and transport along an inverse path.
For a set $S$ and a family $P:S\to\mathrm{Set}$, equality of pairs
$(s,x)=(t,y)$ gives equality $\mv_P(p,x)=y$ over any specified path
$p:s=t$: extract its base path and fibre equation, then use sethood of
$S$ to identify the extracted base path with $p$. In particular pairing
with a fixed $s$ is injective. Conversely, the base path and that fibre
equation give equality of pairs. We use these named statements instead
of treating equality of pairs as equality of their second projections.

\section{The total functor}
Define
\[
 (\mathsf T F)(B)=\sum_{\iota:k\to B}F_B(\iota),\qquad
 (\mathsf T F)(h)(\iota,x)=(h\iota,F(c_{h,\iota})(x)),
\]
where $c_{h,\iota}:(B,\iota)\to(C,h\iota)$ is the canonical algebra
map with underlying map $h$ and reflexive commuting equation.
The values are sets, by the sethood of ring maps and of the fibres.

\begin{lemma}[Change of target structure map]\label{lem:change}
Given $p:h\iota=\kappa$, transporting $F(c_{h,\iota})(x)$ along $p$
gives $F(h,p)(x)$ in $F_C(\kappa)$.
\end{lemma}
\begin{proof}
Induct on $p$. At reflexivity the two algebra maps have the same ring map;
their commuting proofs agree because they are propositions. Extensionality
of algebra maps identifies them, and applying $F$ gives the equation.
\end{proof}

Consequently, for every actual algebra map $a:(B,\iota)\to(C,\kappa)$,
\begin{equation}\label{eq:total-on}
 (\mathsf T F)(\underline a)(\iota,x)=(\kappa,F(a)(x)).
\end{equation}
The base path is the commuting equation of $a$; the fibre equation is
Lemma~\ref{lem:change}. Applying this to the algebra identity proves
the total identity law. For composition, use the two canonical algebra
maps successively and their actual algebra composite. Equation
\eqref{eq:total-on} compares the direct total action with that composite,
and the relative composition law compares the fibre values. Thus the
whole functor law is proved, including its base equation.
Projection to $\iota$ is a natural map $\pi_F:\mathsf T F\to h_k$.

A relative natural transformation totalises by
$(\iota,x)\mapsto(\iota,\alpha_{(B,\iota)}x)$.
Its naturality is relative naturality on canonical algebra maps;
\eqref{eq:total-on} expresses it on arbitrary algebra maps.
This preserves identities and composition and preserves projection to
$h_k$. Relative natural equivalences totalise to natural equivalences,
using their pointwise inverses over each unchanged structure map.

\section{Covers and the sheaf argument}
A finite basic cover of $(B,\iota)$ uses a finite basic cover of $B$.
Put $B_i=B[1/a_i]$, $B_{ij}=B[1/(a_i a_j)]$, with maps
$\lambda_i:B\to B_i$ and $\lambda_{ij}:B\to B_{ij}$.
The relative objects have structure maps $\lambda_i\iota$ and
$\lambda_{ij}\iota$. The overlap maps are the original localisation
maps, supplied with their commuting equations. The two composites from
$(B,\iota)$ to the overlap have the same ring map, hence are equal as
algebra maps. Relative compatibility, matching, and the sheaf condition
are stated on this actual algebra diagram. Empty covers are included.

\begin{lemma}[Fixed-base cover comparison]\label{lem:cover}
Pairing relative values with their specified structure maps identifies
relative compatibility and matching with the corresponding total
compatibility and matching for these paired values.
\end{lemma}
\begin{proof}
Equation~\eqref{eq:total-on} identifies the two reindexings with pairs
whose first component is the same specified overlap structure map.
An equality of relative values gives an equality of those pairs.
Conversely injectivity of pairing at that fixed structure map recovers
the relative equality. Apply this on every overlap or cover member.
\end{proof}

\begin{theorem}[Total sheafhood]\label{thm:sheaf}
A relative functor $F$ is a sheaf if and only if $\mathsf T F$ is an
absolute Zariski sheaf.
\end{theorem}
\begin{proof}
Suppose $F$ is a relative sheaf and $(\iota_i,x_i)$ is a compatible total
family. Projection gives a compatible family of ring maps $k\to B_i$.
The representable sheaf $h_k$ supplies a unique glued map $\iota:k\to B$
and equations $p_i:\lambda_i\iota=\iota_i$.
Rebase $x_i$ along $p_i^{-1}$, obtaining
$y_i\in F(B_i,\lambda_i\iota)$. Path induction gives the equation
$(\lambda_i\iota,y_i)=(\iota_i,x_i)$ in the total value set.
Apply total reindexing to these equations. Together with the original
compatibility they prove compatibility of the paired $y_i$.
Lemma~\ref{lem:cover} gives relative compatibility. Glue in $F$ over
$(B,\iota)$; pairing the result with $\iota$ gives a total amalgam.

For uniqueness, any other total amalgam $(\kappa,z)$ projects to a
base amalgam. Base uniqueness gives $p:\kappa=\iota$. Transport $z$
along $p$; the pair equation identifies the new pair with the original
amalgam point. Lemma~\ref{lem:cover} recovers its relative matching
conditions after the rebasing equations. Relative uniqueness identifies
its fibre value with the constructed one. The total values are therefore
equal. Matching conditions are propositions, so the amalgams themselves
are equal. This proves contractibility of total amalgams.

Conversely, pair a compatible relative family with its prescribed base
maps and glue in the total sheaf. Its projected base map has the same
restrictions as $\iota$; separation of $h_k$ identifies it with $\iota$.
Rebase the total glued value into $F_B(\iota)$. The pair equations and
Lemma~\ref{lem:cover} give relative matching. Two relative amalgams give
total amalgams, so total uniqueness and injectivity of fixed-base pairing
prove relative uniqueness. All these arguments apply to empty covers.
\end{proof}

\section{Compact opens and affine charts}
A relative compact open assigns $U_{(B,\iota)}(x)\in\operatorname{Zar}(B)$
naturally in algebra maps. Define
\[
 (\mathsf T U)_B(\iota,x)=U_{(B,\iota)}(x).
\]
Its naturality is that of $U$ on the canonical map. Conversely an open of
$\mathsf T F$ restricts to a relative open by pairing with $\iota$;
\eqref{eq:total-on} proves naturality on arbitrary algebra maps.
The two operations preserve pointwise covers and inverse one another.
Reordering $(\iota,(x,m))$ as $((\iota,x),m)$ gives a natural equivalence
$\mathsf T|U|\simeq|\mathsf T U|$, with the same point and membership
proof; its naturality uses the same total-action equation.

For a test algebra $(A,\alpha)$, let
$h^k_{(A,\alpha)}(B,\iota)=\Hom_k((A,\alpha),(B,\iota))$.
Forgetting the base map and its commuting equation gives
\begin{equation}\label{eq:affine}
 \mathsf T h^k_{(A,\alpha)}\simeq h_A.
\end{equation}
Indeed the inverse sends $h:A\to B$ to the base map $h\alpha$ and
the canonical algebra map. Its base map is forced by $h$, and the
remaining commuting equation is a proposition. The fibre comparison is
not an identification of the two carriers by computation.

It follows that an actual relative affine atlas totalises to an actual
absolute affine atlas. Affine witnesses merely exist; eliminate each
witness into affineness, and the atlas into existence of an atlas.
Together with Theorem~\ref{thm:sheaf}, a relative functorial qcqs scheme
has absolute qcqs totalisation.

For the converse, an absolute affine chart $e:h_A\simeq|V|$ of
$\mathsf T F$ supplies an algebra structure on $A$: let $\alpha:k\to A$
be the base component of the generic point $e_A(1_A)$.
Naturality gives the base equation
$\pi_F(e_B(h))=h\alpha$ for every $h:A\to B$.
Restricting $e$ to the fibres over $\iota$ therefore gives
$h^k_{(A,\alpha)}\simeq|U|$, where $U$ is the relative open associated
with $V$. The inverse must be rebased using this equation; it does not
follow by discarding the first component of a dependent pair.
Hence absolute affineness of each total chart gives relative affineness,
and absolute qcqs of $\mathsf T F$ implies relative qcqs of $F$.

\section{Fibres and representation over the base}
For an absolute functor $G$ and a natural map $q:G\to h_k$, its fibre
functor is
\[
 \mathsf S(G,q)(B,\iota)=\sum_{z:G(B)}(q_B(z)=\iota).
\]
Reindex $z$ along an algebra map and use naturality of $q$ followed by
its commuting equation to get the new fibre proof. These proofs lie in
a set, so point extensionality proves the functor laws. Forgetting the
base map gives $\mathsf T\mathsf S(G,q)\simeq G$ over $h_k$.
For the other direction,
\[
 \mathsf S(\mathsf T F,\pi_F)(B,\iota)
   =\sum_{\kappa:k\to B}\sum_{x:F_B(\kappa)}(\kappa=\iota)
   \simeq F_B(\iota),
\]
by transporting $x$ along the displayed equality. The inverse pairs
$x$ with $\iota$ and reflexivity. Path induction gives both inverse laws,
and Lemma~\ref{lem:change} gives naturality. Thus the recovery of a
relative functor from its totalisation is a genuine fibre equivalence.
The same argument shows that an absolute natural equivalence preserving
the projections restricts to a relative natural equivalence of fibres.

Let $F$ be relative functorially qcqs. The absolute theorem gives merely
$X$ qcqs and $e:\mathsf T F\simeq h_X$. Compose $e^{-1}$ with $\pi_F$
to get $h_X\to h_k$. The affine comparison $h_k\simeq h_{\Spec k}$
and full faithfulness recover $s:X\to\Spec k$ with its equation on
all affine points. Taking fibres over $\iota$ identifies the fibre of
the projection with
\[
 \Hom_{\Spec k}(\Spec(B,\iota),X).
\]
The defining equation of $s$ and injectivity of the affine comparison
supply the fibre equivalence. The preceding recovery equivalence now
gives a relative natural equivalence from $F$ to this relative geometric
functor. All choices of absolute representation are eliminated only into
mere relative representability at this stage.

Conversely, totalising the relative geometric functor of $(X,s)$ gives
$h_X$: the underlying ordinary point determines its unique structure map
$k\to B$ by the affine comparison. Its canonical lift and uniqueness
are the existing total-point argument for relative full faithfulness.
The absolute geometric atlas theorem and the converse totalisation
criterion above give relative qcqs.

\section{Categorical conclusion and review}
Geometric objects are qcqs spaces with a map to $\Spec k$.
More generally, in the slice of a univalent category $C$ over $b$, an
object is $(x,s:x\to b)$ and an arrow to $(y,t)$ is $(f:x\to y,t f=s)$.
An isomorphism is precisely a base isomorphism $e:x\cong y$ with $te=s$:
the inverse commutes since $s e^{-1}=t e e^{-1}=t$. The commuting
equation is a proposition because the hom types are sets.

Fix $(x,s)$. To show its isomorphism singleton contractible, first use
univalence to contract the base singleton $(y,e:x\cong y)$.
At $(x,1_x)$ the remaining type is $(t:x\to b)\times(t1_x=s)$,
which is contractible: the right-unit law identifies its equation with
$t=s$, a based singleton. Decorating the base identity system by this
unit argument makes the whole structured total space contractible.
Its isomorphisms correspond to actual slice isomorphisms, retaining
the whole underlying arrow. The categorical total-space criterion gives
univalence of the slice. This will be applied to the qcqs category and
its object $\Spec k$.

For native relative functors, the index of the value family is the
FIXED type of actual algebra tests $(B,\iota)$; the dependent transport
in totalisation is not reintroduced at this stage. SetIdentity and
FamilyIdentity give an identity system for this value family. Decorate
it by actions on actual algebra maps and their functor laws. At identity
on the value family, naturality is pointwise equality of the action
maps; function extensionality identifies them. The functor laws are
propositions because values are sets. Reordering the resulting
structural isomorphism data gives exactly native natural equivalences,
with the same natural map. The categorical singleton criterion proves
univalence. Adding the propositional relative qcqs condition preserves
univalence.

The relative functor of points uses the existing BasePoints values,
actions, and equations. Its functor laws on arrows are geometric units
and associativity, followed by point extensionality. The existing
relative Hom--Nat theorem proves full faithfulness on these same native
natural transformations. Relative reconstruction above proves mere
essential surjectivity, with no chosen global representing spaces.
The checked general equivalence theorem then removes truncation only
in the propositional representation fibres and supplies the relative
categorical equivalence with an explicit inverse functor and its two
natural isomorphisms.

The review obligations are: state the transport in $F_B$ before using
it; prove Lemma~\ref{lem:change} before a total functor law; distinguish
an arbitrary algebra map from its canonical target lift; state the
base-gluing equations before moving local values; extract a fibre
equation only over a specified base path; and prove every fibre inverse
law with the base proof still present. The base ring-map type is a set,
so its path witnesses can be identified, but the value fibres vary and
must never be identified by proof irrelevance. Atlas and representation
truncations are eliminated into propositions, and empty covers and the
zero ring are retained. No core or category refactoring is needed.

\begin{theorem}[Relative comparison]
A native small-set-valued functor on small $k$-algebras is functorially
qcqs if and only if there merely exist a small geometric qcqs scheme
$X$, a map $s:X\to\Spec k$, and a native natural equivalence
\[
 F(B,\iota)\simeq
 \Hom_{\Spec k}(\Spec(B,\iota),(X,s)).
\]
The relative functor of points gives an equivalence of univalent
categories
\[
 \mathsf{Qcqs}_\ell/\Spec k\ \simeq\ \mathsf{FunQcqs}_{k,\ell},
\]
with an explicit inverse functor and both natural isomorphisms.
\end{theorem}
\begin{proof}
Sections 3--4 identify intrinsic relative qcqs with absolute qcqs of the
totalisation. Section 5 gives both directions of relative representation.
The slice and native functor-category arguments above give univalence.
Relative full faithfulness and mere essential surjectivity then give
the categorical equivalence, including the stated quasi-inverse.
\end{proof}

\section{Checked correspondence with the detailed argument}
This argument and its review were written before formalisation. The
relative totalisation theorem, both directions of relative essential
image, and the categorical comparison are now checked in Unit. The
proof follows the preceding sections, using the native algebra tests
and the existing geometric maps throughout.

\begingroup\raggedright
Section 1 corresponds to \texttt{DependentPairs}, \texttt{RelativeTests},
and \texttt{RelativeFunctors}. In particular
\texttt{fibre-equality} extracts a fibre equation over a specified base
path, \texttt{rebase-pair} retains the complete pair equation, and
\texttt{push-at} states evaluation of a composite ring map once.

Section 2 corresponds to \texttt{RelativeTotal},
\texttt{RelativeEquivalences}, and \texttt{RelativeTotalMaps}.
\texttt{RelativeTotal.change-target} is Lemma~\ref{lem:change};
\texttt{on-algebra} is equation~\eqref{eq:total-on}. The identity and
composition laws use this equation on actual algebra maps. The total
map laws and its equation over the base are exposed separately.

Section 3 corresponds to \texttt{RelativeCovers},
\texttt{RelativeSheaves}, \texttt{RelativeCoverFibre},
\texttt{RelativeBaseGlue}, \texttt{RelativeTotalSheaf}, and
\texttt{RelativeSheafOfTotal}. The proof names the pointwise base square
before applying algebra-map extensionality. Base gluing, its matching
paths, the rebased local values, and their pair equations are separate
declarations. The two sheaf arguments then use exactly those interfaces.

Section 4 corresponds to \texttt{RelativeRestriction},
\texttt{RelativeOpens}, \texttt{RelativeRealisationTotal},
\texttt{RelativeTotalAffine}, and \texttt{RelativeAffineOfTotal},
followed by \texttt{RelativeQcqs}, \texttt{RelativeTotalQcqs}, and
\texttt{RelativeQcqsOfTotal}.
\texttt{RelativeTotalisation.sheaf-equivalence} and
\texttt{qcqs-equivalence} package the two intrinsic-total criteria.
The converse affine proof names the generic algebra structure and the
base path before transporting the value; its trace retains the total
point and supplies naturality and inverse laws.

Section 5 corresponds to \texttt{RelativeFibres},
\texttt{TotalFibreEquivalence}, \texttt{RelativeFibreRecovery},
\texttt{RelativeFibreEquivalence}, \texttt{RelativeGeometric},
\texttt{RelativeGeometricFibre}, and \texttt{GeometricBaseRecovery}.
\texttt{RelativeQcqsTransfer} and \texttt{RelativeGeometricQcqs}
give the invariance and geometric example needed for
\texttt{RelativeEssentialImage.reconstruct} and \texttt{recognise}.
The recovered structure map is opaque after its equation on affine
points. The geometric fibre inverse is opaque after its point equation,
so its inverse-law clients do not reopen its base-equation proof.

Section 6 corresponds to \texttt{SliceCategory},
\texttt{RelativeFunctorIdentity}, \texttt{RelativeFunctorCategory},
and \texttt{RelativeCategories}. The final theorem is
\texttt{RelativeCategories.Over[K].equivalence}; the object equivalence
is also exposed separately. Its inverse functor and both natural
isomorphisms are the checked general construction of
\texttt{CategoryQuasiInverse}.

Validation: 249 modules, 9,250 declarations, zero failures. The combined
\texttt{essential-image.opens}, \texttt{categories.opens}, and
\texttt{relative.opens} audit reports no extra openings in any of its
3,752 lists. The 35 new modules have 863 audited declarations; each takes
at most 0.05 seconds in the isolated audit, and at most 33 distinct
names are opened. The construction proofs of function extensionality,
quotients, and localisation are not reopened by the new clients.
The comparison adds no axiom, choice, or resizing principle and changes
no core or category-library implementation.
\par\endgroup
\end{document}
